% Imporditud: tools/import_eksperimendid.py
% Allikas: Eesti füüsikaolümpiaadi eksperimendi ülesannete
%   kogu (Rael Kalda, 2023), ülesanne Ü133, lahendus L133.
\setAuthor{EFO žürii}
\setRound{lõppvoor}
\setYear{2007}
\setNumber{G E2}
\setDifficulty{4}
\setTopic{elektriahelad}
\setVahendid{must kast punase, rohelise ja kollase väljundjuhtmega, milles on teatud viisil ühendatud kolm komponenti (neist kaks on ühesugused), patarei, lamp, mille värv sõltub voolusuunast}
\setPunktid{14}
\setAllikas{Eksperimendi kogu Ü133, lahendus lk 102}
\setLahendusseis{puudub}

\prob{Must kast}
Määrata, mis tüüpi komponendid on mustas kastis ja teha kindlaks musta kasti elektriline skeem. Kirjeldage, milliste katsete põhjal ja kuidas jõudsite oma järeldusele.

\hint

\solu
\textit{Kogumikus lahendus puudub.}
\probend
